\section{Exercises: Finding Needs}

\subsection{From Feelings to Needs}

\begin{enumerate}
    \item \textbf{``I feel very lonely, as if I am not important.'' --- What is the need behind this?}
    \begin{itemize}
        \item \textbf{Analysis:} Loneliness usually points to unmet needs
        for connection, companionship, and belonging; ``not important''
        points to the need to be seen and valued.
        \item \textbf{Need:} Connection, companionship.
        \item \textbf{Expression Practice:} ``I feel very lonely, because I
        deeply desire to feel that we are authentically connected.''
    \end{itemize}

    \item \textbf{``I get annoyed as soon as I hear her voice.'' --- What is the need behind this?}
    \begin{itemize}
        \item \textbf{Analysis:} Annoyance is often a manifestation of
        elevated cortisol, indicating a core need has been unmet for a long
        time \citep{howard2024}. Here we need to ask: In this relationship,
        what has not been satisfied?
        \item \textbf{Possible Needs:} Respect, autonomy, tranquility,
        connection.
        \item \textbf{Expression Practice:} ``Every time we talk, I feel
        very annoyed, I think it's because I deeply desire to feel equally
        treated in our communication.''
    \end{itemize}

    \item \textbf{``He never thanks me for everything I do.'' --- What is the need behind this?}
    \begin{itemize}
        \item \textbf{Analysis:} ``Never'' is an absolute judgment, hiding
        long-suppressed grievances. The real need is for contribution to
        be seen \citep{izard2001}.
        \item \textbf{Need:} Understanding, fairness.
        \item \textbf{Expression Practice:} ``I feel somewhat lost, because
        I long for my efforts to be understood and treated fairly.''
    \end{itemize}

    \item \textbf{``I am really worried about our future.'' --- What is the need behind this?}
    \begin{itemize}
        \item \textbf{Analysis:} Worry is a signal of rising cortisol,
        warning that a core need involving the future is facing a threat \citep{herman2025}.
        \item \textbf{Need:} Security, clarity, and consensus.
        \item \textbf{Expression Practice:} ``I feel somewhat worried,
        because I deeply desire that we can jointly have a clear direction
        and consensus about our future.''
    \end{itemize}

    \item \textbf{``Every time I give an opinion, he says I'm being nosy.'' --- What is the need behind this?}
    \begin{itemize}
        \item \textbf{Analysis:} Being denied the right to participate
        hides the need for contribution and acceptance \citep{hoffman2000}.
        \item \textbf{Need:} Respect.
        \item \textbf{Expression Practice:} ``Every time I share my
        thoughts, I feel somewhat hurt, because I deeply desire my thoughts
        to be respected.''
    \end{itemize}

    \item \textbf{``I don't know why, but I can't get motivated to do anything.'' --- What is the need behind this?}
    \begin{itemize}
        \item \textbf{Analysis:} This is a typical description of low
        dopamine and a lack of meaning \citep{dweck2006}.
        \item \textbf{Need:} Meaning, purpose, creation, connection, growth.
        \item \textbf{Expression Practice:} ``Recently I feel somewhat dull,
        I think I internally desire to do some things that make me feel
        truly meaningful.''
    \end{itemize}

    \item \textbf{``He always turns a deaf ear to my words.'' --- What is the need behind this?}
    \begin{itemize}
        \item \textbf{Analysis:} ``Always'' is a judgment, hiding the unmet
        needs of being listened to and valued \citep{rogers1961}.
        \item \textbf{Need:} Connection, respect.
        \item \textbf{Expression Practice:} ``Every time I finish speaking
        and get no response, I feel somewhat lost, because I deeply desire
        to feel connected.''
    \end{itemize}

    \item \textbf{``I feel that he is controlling everything about me.'' --- What is the need behind this?}
    \begin{itemize}
        \item \textbf{Analysis:} ``Being controlled'' is a pseudo-feeling,
        hiding a severe lack of the needs for autonomy and freedom \citep{thomas1977}.
        \item \textbf{Need:} Autonomy.
        \item \textbf{Expression Practice:} ``When my schedule and choices
        need to pass your approval, I feel very oppressed, because I deeply
        desire to be able to freely make my own choices.''
    \end{itemize}

    \item \textbf{``Just thinking about going to work tomorrow makes me feel bad.'' --- What is the need behind this?}
    \begin{itemize}
        \item \textbf{Analysis:} This continuous resistance often comes
        from the systemic lack of a core need in the work environment \citep{lazarus1991}.
        \item \textbf{Need:} Meaning, respect, autonomy, belonging, rest \citep{kabatzinn1990}.
        \item \textbf{Expression Practice:} ``Just thinking about tomorrow,
        I feel somewhat heavy. I think what I internally desire is an
        environment that makes me feel my work is meaningful.''
    \end{itemize}
\end{enumerate}

\begin{keyinsight}
Identifying the underlying need behind our own and others' surface feelings
requires intentional practice.
Whether translating anger into a desire for respect or sadness into a plea
for connection, mastering this skill is essential for breaking negative
communication cycles.
\end{keyinsight}
